\section{Introduction}

Take the ring $\mathbb C[p_1,p_2,\ldots]$ of polynomials in countably many variables $p_1, p_2,\ldots$ and
let $\alpha\in\mathbb C$ be a~parameter.
For $n=1,2,\ldots$ denote by $p_n^{\ast}$ the operator
$\alpha n\partial/\partial p_n$ of renormalized dif\/ferentiation relative to the
variable $p_n$.
The following dif\/ferential operator on $\mathbb C[p_1,p_2,\ldots]$ arises in various problems of
algebra and mathematical physics and has recently enjoyed considerable attention:
\begin{gather}
\label{def-I}
\sum_{m,n\ge1}\big(p_mp_np_{m+n}^{\ast}+p_{m+n}p_m^{\ast}p_n^{\ast}\big)+(\alpha-1)\sum_{n\ge1}np_np_n^{\ast},
\end{gather}
see~\cite{AMOS,CJ,NS1,OP,Pol} for instance.
In the classical limit $\alpha\rightarrow0$, the parameter $\alpha$ playing the role of the Planck
constant, the operator~\eqref{def-I} becomes the Hamiltonian of the periodic Benjamin--Ono equation,
a~well-studied integrable Hamiltonian system~\cite{AFSS,B,KLM,KM,O}.

The ring $\mathbb C[p_1,p_2,\ldots]$ can be identif\/ied with the ring of symmetric functions in the
variables $x_1,x_2,\ldots$ by setting $p_n=x_1^{ n}+x_2^{ n}+\cdots$ for
every $n$.
It is known that then the operator~\eqref{def-I} is diagonalised in the basis of the Jack symmetric
functions corresponding to the parameter~$\alpha$.
It is also known~\cite{SV} that~\eqref{def-I} can be included into an inf\/inite family of commuting
operators, all diagonal in the basis of Jack symmetric functions.
In~\cite{NS1} we provided an explicit expression for a~generating function $A(u)$ of those commuting
operators, see Theorem~\ref{Theorem1} in Section~\ref{Section4} of the present article.
The function $A(u)$ has been def\/ined as the projective limit of the renormalized Sekiguchi--Debiard
determinant for the $N$-particle Calogero--Sutherland model as $N\rightarrow\infty$.

The purpose of the present article is to construct a~generating function $I(u)$ of another family
$I^{(1)},I^{(2)},\ldots$ of commuting operators diagonal in the basis of
the Jack symmetric functions.
Our construction is drastically dif\/ferent from that of~\cite{NS1} and has been inspired by the Lax
formulation of the classical Benjamin--Ono equation.
Namely, put
\begin{gather}
\label{iu}
I(u)=I^{(1)}u^{-1}+I^{(2)}u^{-2}+\cdots,
\end{gather}
where by def\/inition
\begin{gather}
\label{hdef}
1-u^{-1}I(u)=A(u)/A(u+1).
\end{gather}
Our main result is Theorem~\ref{Theorem2} in Section~\ref{Section4} giving an explicit formula for $I(u)$.
In particular $I^{(1)}$ is the weight counting operator
$p_1p_1^{\ast}+p_2p_2^{\ast}+\cdots$ while $I^{(2)}$
coincides with the operator~\eqref{def-I}.

An advantage of this new generating function $I(u)$ is that it is written explicitly in terms of the
collective variables $p_1,p_2,\ldots$ whereas $A(u)$ is written in terms of the monomial symmetric
functions, which are expressed in these variables in a~rather complicated combinatorial way.

